\documentclass[UTF8,12pt]{ctexart}
\usepackage[a4paper,margin=24mm]{geometry}
\usepackage{amsmath,amssymb,bm,graphicx,adjustbox}
\newcommand{\fitmath}[1]{\begin{adjustbox}{max width=\linewidth}$\displaystyle #1$\end{adjustbox}}
\setlength{\parindent}{0pt}
\setlength{\parskip}{.7em}
\sloppy
\allowdisplaybreaks
\begin{document}
\section*{1988 年考研数学三 · 参考答案}
\subsection*{1988年真题参考答案}

\subsection*{（试卷IV）}

\subsection*{一、填空题}

（1）\textcircled{1} \(\displaystyle e^{-x^2/2}\)；\textcircled{2} 在 \(\displaystyle (-\infty,\allowbreak +\infty)\) 上单调增；\textcircled{3} 奇函数；\textcircled{4} \(\displaystyle (0,\allowbreak 0)\)；\textcircled{5} 在 \(\displaystyle (-\infty,\allowbreak 0)\) 上是凹的，在 \(\displaystyle (0,\allowbreak +\infty)\) 上是凸的；\textcircled{6} \(\displaystyle y=-\sqrt{\pi/2},\allowbreak \ y=\sqrt{\pi/2}\)。


（2）\(\displaystyle -3\)。


（3）


\[\fitmath{\displaystyle \begin{pmatrix}0&0&0&1\\0&0&1&0\\0&1&0&0\\1&0&0&0\end{pmatrix}.}\]


（4）\textcircled{1} \(\displaystyle 0.3\)；\textcircled{2} \(\displaystyle 0.5\)。


\subsection*{二、判断题}

（1）×。 （2）×。 （3）×。 （4）√。 （5）×。


\subsection*{三、计算题}

（1）\(\displaystyle 1\)。 （2）\(\displaystyle \dfrac{\displaystyle \partial^2u}{\displaystyle \partial x\partial y}=\dfrac{\displaystyle 1}{\displaystyle 1+e^u}-\dfrac{\displaystyle xye^u}{\displaystyle (1+e^u)^3}\)。 （3）\(\displaystyle \dfrac{\displaystyle 2\pi}{\displaystyle 3}\)。 （4）\(\displaystyle \dfrac{\displaystyle 1}{\displaystyle 2}\)。


\subsection*{四、级数}

（1）级数收敛。（2）证明略（可利用 \(\displaystyle |a_nb_n|\le\dfrac{\displaystyle 1}{\displaystyle 2}(a_n^2+b_n^2)\) 和比较审敛法）。


五、（1）\(\displaystyle p_e=(a/b)^{1/3}\)。 （2）\(\displaystyle p(t)=[p_e^3+(1-p_e^3)e^{-3kbt}]^{1/3}\)。 （3）\(\displaystyle \lim_{t\to+\infty}p(t)=p_e\)。


六、（1）\(\displaystyle (1,\allowbreak 1)\)。 （2）\(\displaystyle y=2x-1\)。 （3）\(\displaystyle \dfrac{\displaystyle \pi}{\displaystyle 30}\)。


七、\(\displaystyle k_1\ne2\) 时，方程组有唯一解；\(\displaystyle k_1=2\) 且 \(\displaystyle k_2\ne1\) 时，方程组无解；\(\displaystyle k_1=2\) 且 \(\displaystyle k_2=1\) 时，方程组有无穷多解，其解为 \(\displaystyle (x_1,\allowbreak x_2,\allowbreak x_3,\allowbreak x_4)^{\mathrm T}=k(0,\allowbreak -2,\allowbreak 1,\allowbreak 0)^{\mathrm T}+(-8,\allowbreak 3,\allowbreak 0,\allowbreak 2)^{\mathrm T}\)，其中 \(\displaystyle k\) 为任意常数。


八、\(\displaystyle s\) 为奇数时，向量组 \(\displaystyle \boldsymbol\beta_1,\allowbreak \ldots,\allowbreak \boldsymbol\beta_s\) 线性无关；\(\displaystyle s\) 为偶数时，向量组 \(\displaystyle \boldsymbol\beta_1,\allowbreak \ldots,\allowbreak \boldsymbol\beta_s\) 线性相关。


九、\(\displaystyle -\dfrac{\displaystyle 16}{\displaystyle 27}\)。


十、（1）\(\displaystyle \alpha=\dfrac{\displaystyle 448}{\displaystyle 475}\approx0.94\)。 （2）\(\displaystyle \beta\approx0.85\)。


十一、（1）\(\displaystyle P\{X=k\}=\dbinom{\displaystyle 100}{\displaystyle k}0.2^k0.8^{100-k}\;(k=0,\allowbreak 1,\allowbreak \ldots,\allowbreak 100)\)。 （2）\(\displaystyle 0.927\)。


十二、


\[\fitmath{\displaystyle f(y)=\begin{cases}\dfrac{\displaystyle 1}{\displaystyle 2y},&e^2<y<e^4,\\0,&\text{其他}.\end{cases}}\]


\subsection*{（试卷V）}

一、【同试卷IV 第一题】 二、【同试卷IV 第二题】


三、（1）\(\displaystyle \dfrac{\displaystyle 4}{\displaystyle \pi}\)。 （2）\(\displaystyle -\dfrac{\displaystyle x+y}{\displaystyle y^3}e^{x/y}\)。


三、（3）【同试卷IV 第三、（3）题】 （4）【同试卷IV 第三、（4）题】


四、\(\displaystyle a=2,\allowbreak b=-1\)。


五、当圆的周长为 \(\displaystyle \dfrac{\displaystyle \pi a}{\displaystyle 4+\pi}\)，正方形的周长为 \(\displaystyle \dfrac{\displaystyle 4a}{\displaystyle 4+\pi}\) 时，两图形的面积之和最小。


六、【同试卷IV 第六题】 七、【同试卷IV 第七题】


八、证明略。\(\displaystyle A^{-1}=\dfrac{\displaystyle 1}{\displaystyle 2}(A-3E)\)。


九、【同试卷IV 第八题】 十、【同试卷IV 第十题】


十一、（1）\(\displaystyle P\{X=0\}=\dfrac{\displaystyle 4}{\displaystyle 5},\allowbreak \ P\{X=1\}=\dfrac{\displaystyle 8}{\displaystyle 45},\allowbreak \ P\{X=2\}=\dfrac{\displaystyle 1}{\displaystyle 45}\)。


（2）\(\displaystyle E(X)=\dfrac{\displaystyle 2}{\displaystyle 9}\)。 （3）\(\displaystyle D(X)=\dfrac{\displaystyle 88}{\displaystyle 405}\)。


十二、【同试卷IV 第十二题】

\end{document}
